\subsection{Native estimators}\label{sec:estimators_python}

This section describes the factory functions provided by
\gretlpy for accessing native \lgretl estimators. Each
function creates and returns an instance of the corresponding
model-specific class, derived from the abstract
\classpymodel base class. These functions do not perform
estimation themselves; instead, they construct a model
specification that can subsequently be estimated by calling
the \pythcode{fit()} method.

All estimator functions share a common calling convention.
The \textit{formula} argument specifies the dependent
variable, regressors, and, where applicable, additional
elements required by a particular estimator. Depending on
the function, the specification may be supplied as a string
or a Python list. The supported forms of model specification
are discussed in Section~\ref{sec:estimators_python_examples}.

Estimator-specific options are passed as keyword arguments
(\textit{OPTS}). These generally correspond to the options
available for the respective commands in \lgretl. In addition,
all estimator functions accept the optional keyword argument
\pythcode{data}, which takes an instance of
\classpydataset and specifies the dataset associated with
the model. If \pythcode{data} is omitted, the currently
default dataset is used (see \ref{sec:globalandindivodual()}
for details).

The descriptions below provide the Python function
signatures, their correspondence to native \lgretl commands,
and the supported estimation options. The return type
indicates the model-specific class created by each factory
function. Estimation results become available through the
returned object after calling \pythcode{fit()}.

\gretlfpycmd{ar()}
\estdesc{Specify a linear regression model with autoregressive
	errors, estimated using the generalized Cochrane--Orcutt
	iterative procedure.}
\begin{lstlisting}[style=pythonstylefoot]
gretl.ar(formula: str | list[, OPTS]) -> gretl.Model_ar
\end{lstlisting}
\mapto[none]{ar ~lags ; depvar indepvars}
\noRformula
\begin{itemize}
	\item[]\maptoopt{vcv=True}{--vcv}{print covariance matrix}
	\item[]\maptoopt{quiet=True}{--quiet}{suppress printing of parameter estimates}
\end{itemize}
\generaldesc{The model specification begins with a list of
	residual lags, followed by a semicolon and the dependent
	variable and regressors. The residual lags need not be
	consecutive.
	
	When using Python lists, the residual lags and model variables
	are represented by nested lists. For example, the specifications
	\pythcode{"1 3 4 ; y 0 x1 x2 x3"} and
	\pythcode{[[1, 3, 4], ["y", 0, "x1", "x2", "x3"]]}
	are equivalent.
	
	Estimation uses the generalized Cochrane--Orcutt iterative
	procedure. Iteration terminates when successive error sums
	of squares differ by no more than 0.005 percent, or after
	20 iterations.}

\gretlfpycmd{ar1()}
\estdesc{Specify a linear regression model with first-order
	autoregressive errors, estimated by feasible generalized
	least squares.}
\begin{lstlisting}[style=pythonstylefoot]
gretl.ar1(formula: str | list[, OPTS]) -> gretl.Model_ar1
\end{lstlisting}
\mapto{ar1 depvar indepvars}
\begin{itemize}
	\item[]\maptoopt{hilu=True}{--hilu}{use Hildreth--Lu procedure}
	\item[]\maptoopt{pwe=True}{--pwe}{use Prais--Winsten estimator}
	\item[]\maptoopt{vcv=True}{--vcv}{print covariance matrix}
	\item[]\maptoopt{no\_corc=True}{--no-corc}{skip Cochrane--Orcutt refinement of Hildreth--Lu estimates}
	\item[]\maptoopt{loose=True}{--loose}{use a looser convergence criterion}
	\item[]\maptoopt{quiet=True}{--quiet}{suppress printing of results}
\end{itemize}
\generaldesc{By default, estimation uses the iterative
	Cochrane--Orcutt procedure. Convergence is achieved when
	successive estimates of the autocorrelation coefficient
	differ by no more than $10^{-6}$, or by no more than
	$0.001$ if \pythcode{loose=True} is specified.
	An error is reported if convergence is not achieved
	within 100 iterations.
	
	With \pythcode{pwe=True}, the Prais--Winsten estimator
	is used. Unlike Cochrane--Orcutt, this procedure retains
	the first observation.
	
	With \pythcode{hilu=True}, the Hildreth--Lu search
	procedure is used. Its results are subsequently refined
	using Cochrane--Orcutt iterations, unless
	\pythcode{no\_corc=True} is specified. This option has
	no effect with other estimators.}

\gretlfpycmd{arima()}
\estdesc{Specify an ARIMA, SARIMA or ARIMAX model for
	maximum-likelihood estimation.}
\begin{lstlisting}[style=pythonstylefoot]
gretl.arima(depvar: str, order: tuple[, seasonal: tuple = None, indepvars: str | list = None, OPTS]) -> gretl.Model_arima
\end{lstlisting}
\mapto{arima p d q [; P D Q] ; depvar [indepvars]}
\begin{itemize}
	\item[]\maptoopt{verbose=True}{--verbose}{print details of numerical iterations}
	\item[]\maptoopt{quiet=True}{--quiet}{suppress printing of results}
	\item[]\maptoopt{vcv=True}{--vcv}{print covariance matrix}
	\item[]\maptoopt{hessian=True}{--hessian}{disable fallback to OPG}
	\item[]\maptoopt{opg=True}{--opg}{use the outer product of the gradient}
	\item[]\maptoopt{nc=True}{--nc}{do not include a constant}
	\item[]\maptoopt{conditional=True}{--conditional}{use conditional maximum likelihood}
	\item[]\maptoopt{x\_12\_arima=True}{--x-12-arima}{use X-12-ARIMA or X-13ARIMA-SEATS}
	\item[]\maptoopt{lbfgs=True}{--lbfgs}{use the L-BFGS-B maximizer}
	\item[]\maptoopt{y\_diff\_only=True}{--y-diff-only}{difference only the dependent variable}
	%\item[]\maptoopt{lagselect=True}{--lagselect}{perform lag-order selection}
\end{itemize}
\generaldesc{The argument \pythcode{depvar} specifies the
	dependent variable, while \pythcode{order} defines the
	non-seasonal ARIMA orders $(p,d,q)$. The optional
	\pythcode{seasonal} argument specifies the seasonal
	orders $(P,D,Q)$. Exogenous regressors can be supplied
	via \pythcode{indepvars} to specify an ARIMAX model.
	
	In the absence of exogenous regressors, a constant is
	included by default unless \pythcode{nc=True} is specified.
	When exogenous regressors are present, the constant must
	be included explicitly in \pythcode{indepvars}.
	
	By default, estimation uses native exact maximum likelihood.
	Conditional maximum likelihood can be selected using
	\pythcode{conditional=True}. Alternatively, estimation
	can be performed using X-12-ARIMA or X-13ARIMA-SEATS.
	
	Under native exact ML, the covariance matrix is based
	on the numerical Hessian, with fallback to the outer
	product of the gradient (OPG). The mutually exclusive
	options \pythcode{hessian=True} and \pythcode{opg=True}
	control this behavior. The \pythcode{lbfgs=True} option
	selects the L-BFGS-B maximizer instead of BFGS.
	
	For ARIMAX models with non-zero differencing, both the
	dependent variable and the regressors are differenced
	by default. The \pythcode{y\_diff\_only=True} option
	leaves the regressors in levels.
	
	%With \pythcode{lagselect=True}, the AR and MA orders are
	%treated as maximum orders, and a table of information
	%criteria and log-likelihood values is produced instead
	%of estimates for a single specification.
}

\gretlfpycmd{biprobit()}
\estdesc{Specify a bivariate probit model for estimation
	by maximum likelihood, allowing the two equations to
	share regressors or have separate sets of regressors.}
\begin{lstlisting}[style=pythonstylefoot]
gretl.biprobit(formula: str | list[, OPTS]) -> gretl.Model_biprobit
\end{lstlisting}
\mapto{biprobit depvar1 depvar2 indepvars1 [; indepvars2]}
\begin{itemize}
	\item[]\maptoopt{vcv=True}{--vcv}{print covariance matrix}
	\item[]\maptoopt{robust=True}{--robust}{use QML robust standard errors}
	\item[]\maptoopt{cluster="clustvar"}{--cluster=clustvar}{compute clustered standard errors}
	\item[]\maptoopt{opg=True}{--opg}{use the Outer Product of the Gradient covariance matrix}
	\item[]\maptoopt{save\_xbeta=True}{--save-xbeta}{save the index functions as fitted values}
	\item[]\maptoopt{verbose=True}{--verbose}{print additional information}
\end{itemize}
\noRformula\\
\generaldesc{The first two variables specify the binary dependent
	variables. If only one set of regressors is provided, it is used
	in both equations. Otherwise, the regressors preceding the
	semicolon apply to the first equation, and those following
	the semicolon apply to the second equation.
	
	When using Python lists, semicolon-separated variable lists
	are represented by nested lists. For example, the specifications
	\pythcode{"y1 y2 0 x1 x2 ; 0 x3 x4"} and
	\pythcode{[["y1", "y2", 0, "x1", "x2"], [0, "x3", "x4"]]}
	are equivalent.
	
	By default, standard errors are computed using the analytical
	Hessian. The \pythcode{opg} option selects the Outer Product
	of the Gradient covariance matrix, while \pythcode{robust}
	requests QML sandwich standard errors. The \pythcode{cluster}
	option computes cluster-robust standard errors.
	
	The estimated correlation parameter $\rho$ is included as
	the last element of the coefficient vector.
	
	By default, fitted values contain the probabilities of the
	four joint outcomes $(1,1)$, $(1,0)$, $(0,1)$ and $(0,0)$.
	With \pythcode{save\_xbeta=True}, the fitted values contain
	the index functions for the two equations instead.}

\gretlfpycmd{dpanel()}
\estdesc{Specify a dynamic panel-data model for estimation using
	difference GMM (GMM-DIF) or system GMM (GMM-SYS).}
\begin{lstlisting}[style=pythonstylefoot]
gretl.dpanel(p: int | numpy.ndarray, formula: str[, OPTS]) -> gretl.Model_dpanel
\end{lstlisting}
\mapto{dpanel p ; depvar indepvars [ ; instruments ]}
\begin{itemize}
	\item[]\maptoopt{quiet=True}{--quiet}{suppress printing of results}
	\item[]\maptoopt{vcv=True}{--vcv}{print covariance matrix}
	\item[]\maptoopt{two\_step=True}{--two-step}{perform two-step GMM estimation}
	\item[]\maptoopt{system=True}{--system}{add equations in levels (GMM-SYS)}
	\item[]\maptoopt{collapse=True}{--collapse}{collapse GMM-style instruments}
	\item[]\maptoopt{time\_dummies=True}{--time-dummies}{add time dummy variables}
	\item[]\maptoopt{dpdstyle=True}{--dpdstyle}{use the DPD convention for time dummies}
	\item[]\maptoopt{asymptotic=True}{--asymptotic}{use uncorrected asymptotic standard errors}
	\item[]\maptoopt{keep\_extra=True}{--keep-extra}{retain GMM weight and instrument matrices}
\end{itemize}
\noRformula\\
\generaldesc{The argument \textit{p} specifies the autoregressive
	order of the dependent variable. It may be an integer or an array
	of possibly non-contiguous lag orders. The dependent variable and
	regressors are specified in levels; the required differencing is
	performed automatically.
	
	By default, difference GMM (GMM-DIF) is used. Setting
	\pythcode{system=True} adds equations in levels, producing
	the system GMM (GMM-SYS) estimator.
	
	The optional instrument specification follows a semicolon.
	If omitted, all independent variables are assumed to be strictly
	exogenous. Otherwise, strictly exogenous regressors must also
	be included in the instrument list. The \pythcode{GMM()}
	and \pythcode{GMMlevel()} terms specify lagged GMM-style
	instruments for the differenced and levels equations, respectively.
	The \pythcode{collapse} option limits instrument proliferation.
	
	One-step GMM with robust standard errors is used by default.
	With \pythcode{two\_step=True}, the Windmeijer finite-sample
	correction is applied to the standard errors, unless
	\pythcode{asymptotic=True} is specified.
	
	Estimation results include tests for first- and second-order
	autocorrelation, tests of overidentifying restrictions, and
	a Wald test for the joint significance of the regressors.
	First-order autocorrelation in differenced residuals is not
	in itself problematic, whereas second-order autocorrelation
	violates the maintained assumptions.
	
	The \pythcode{time\_dummies} option adds time dummy variables,
	entered in differenced form by default. With
	\pythcode{dpdstyle=True}, they are entered in levels.
	The \pythcode{keep\_extra} option retains the GMM weight
	and instrument matrices in the model results.
	See also \seegretlmanual.}

\gretlfpycmd{duration()}
\estdesc{By default the Weibull distribution is used
	but the exponential, log-logistic and log-normal
	distributions are also available.}
\begin{lstlisting}[style=pythonstylefoot]
gretl.duration(formula: str[, data: gretl.Dataset, OPTS]) -> gretl.Model_duration
\end{lstlisting}
\mapto{duration depvar indepvars [; censvar ]}
\begin{itemize}
	\item[]\maptoopt{exponential=True}{--exponential}{use exponential distribution}
	\item[]\maptoopt{loglogistic=True}{--loglogistic}{use log-logistic distribution}
	\item[]\maptoopt{lognormal=True}{--lognormal}{use log-normal distribution}
	\item[]\maptoopt{medians=True}{--medians}{fitted values are medians}
	\item[]\maptoopt{robust=True}{--robust}{robust standard errors}
	\item[]\maptoopt{vcv=True}{--vcv}{print covariance matrix}
	\item[]\maptoopt{verbose=True}{--verbose}{print details of iterations}
\end{itemize}

\gretlfpycmd{duration()}
\estdesc{Specify a parametric duration model for estimation
	by maximum likelihood. The dependent variable represents
	duration and must be strictly positive. By default, the
	Weibull distribution is used, but exponential, log-logistic
	and log-normal specifications are also available.}
\begin{lstlisting}[style=pythonstylefoot]
gretl.duration(formula: str | list[, OPTS]) -> gretl.Model_duration
\end{lstlisting}
\mapto{duration depvar indepvars [; censvar]}
\begin{itemize}
	\item[]\maptoopt{exponential=True}{--exponential}{use exponential distribution}
	\item[]\maptoopt{loglogistic=True}{--loglogistic}{use log-logistic distribution}
	\item[]\maptoopt{lognormal=True}{--lognormal}{use log-normal distribution}
	\item[]\maptoopt{medians=True}{--medians}{use conditional medians as fitted values}
	\item[]\maptoopt{robust=True}{--robust}{use QML robust standard errors}
	\item[]\maptoopt{cluster="clustvar"}{--cluster=clustvar}{compute clustered standard errors}
	\item[]\maptoopt{vcv=True}{--vcv}{print covariance matrix}
	\item[]\maptoopt{verbose=True}{--verbose}{print details of numerical iterations}
	\item[]\maptoopt{quiet=True}{--quiet}{suppress printing of results}
\end{itemize}
\generaldesc{The dependent variable represents the duration
	of a state or event, such as the length of an unemployment
	spell. All duration values must be strictly positive.
	
	An optional right-censoring variable may be specified after
	a semicolon. Non-zero values of this variable identify
	right-censored observations.
	
	When using Python lists, semicolon-separated variable lists
	are represented by nested lists. For example, the specifications
	\pythcode{"y 0 x1 x2 ; cens"} and
	\pythcode{[["y", 0, "x1", "x2"], ["cens"]]}
	are equivalent.
	
	By default, the fitted values obtained through the
	\gretlcmd{\$yhat} accessor are conditional means of duration.
	With \pythcode{medians=True}, the fitted values are conditional
	medians instead.}

\gretlfpycmd{garch()}
\estdesc{Specify a Generalized Autoregressive Conditional
	Heteroskedasticity (GARCH) model for estimation by maximum
	likelihood. The model may include exogenous regressors
	in the conditional mean equation.}
\begin{lstlisting}[style=pythonstylefoot]
gretl.garch(formula: str | list[, OPTS]) -> gretl.Model_garch
\end{lstlisting}
\mapto[none]{garch p q ; depvar [indepvars]}
\begin{itemize}
	\item[]\maptoopt{robust=True}{--robust}{use QML robust standard errors}
	\item[]\maptoopt{verbose=True}{--verbose}{print details of numerical iterations}
	\item[]\maptoopt{quiet=True}{--quiet}{suppress printing of results}
	\item[]\maptoopt{vcv=True}{--vcv}{print covariance matrix}
	\item[]\maptoopt{nc=True}{--nc}{exclude the constant from the mean equation}
	\item[]\maptoopt{stdresid=True}{--stdresid}{standardize the residuals}
	\item[]\maptoopt{fcp=True}{--fcp}{use the Fiorentini--Calzolari--Panattoni algorithm}
	\item[]\maptoopt{arma\_init=True}{--arma-init}{initialize variance parameters using ARMA}
\end{itemize}
\noRformula\\\\
\generaldesc{The model specification begins with the GARCH
	and ARCH orders $p$ and $q$, followed by a semicolon and
	the dependent variable with optional regressors.
	The order $p$ specifies the GARCH component, while $q$
	specifies the ARCH component. If $p>0$, then $q$ must
	also be positive. The sum $p+q$ must not exceed 5.
	
	A constant is automatically included in the conditional
	mean equation unless \pythcode{nc=True} is specified.
	
	When using Python lists, the lag orders and model variables
	are represented by nested lists. For example, the specifications
	\pythcode{"1 1 ; y x1 x2"} and
	\pythcode{[[1, 1], ["y", "x1", "x2"]]}
	are equivalent.
	
	By default, estimation uses the native BFGS algorithm,
	and the covariance matrix is based on the Hessian.
	The \pythcode{robust} option requests the QML (White)
	covariance matrix, while \pythcode{fcp} selects the
	Fiorentini--Calzolari--Panattoni algorithm.
	
	The variance parameters are initialized using an initial
	OLS regression. With \pythcode{arma\_init=True}, starting
	values are obtained from an ARMA model, which may improve
	convergence.
	
	The residuals and estimated conditional variance are
	available through the \gretlcmd{\$uhat} and
	\gretlcmd{\$h} accessors, respectively.
	With \pythcode{stdresid=True}, the residuals are divided
	by the estimated conditional standard deviation.}

\gretlfpycmd{heckit()}
\estdesc{Specify a Heckman sample-selection model consisting
	of an outcome equation and a selection equation. The dependent
	variable in the selection equation must be binary.}
\begin{lstlisting}[style=pythonstylefoot]
gretl.heckit(formula: str | list[, OPTS]) -> gretl.Model_heckit
\end{lstlisting}
\mapto{heckit depvar indepvars ; selection ~equation}
\begin{itemize}
	\item[]\maptoopt{quiet=True}{--quiet}{suppress printing of results}
	\item[]\maptoopt{two\_step=True}{--two-step}{perform two-step estimation}
	\item[]\maptoopt{vcv=True}{--vcv}{print covariance matrix}
	\item[]\maptoopt{opg=True}{--opg}{use the Outer Product of the Gradient covariance matrix}
	\item[]\maptoopt{robust=True}{--robust}{use QML robust standard errors}
	\item[]\maptoopt{cluster="clustvar"}{--cluster=clustvar}{compute clustered standard errors}
	\item[]\maptoopt{verbose=True}{--verbose}{print additional information}
\end{itemize}
\noRformula\\
\generaldesc{The specification consists of two equations separated
	by a semicolon. The first equation describes the outcome variable,
	while the second specifies the selection mechanism. The dependent
	variable in the selection equation must be binary.
	
	When using Python lists, the two equations are represented
	by nested lists. For example, the specifications
	\pythcode{"y 0 x1 x2 ; ys 0 x3 x4"} and
	\pythcode{[["y", 0, "x1", "x2"], ["ys", 0, "x3", "x4"]]}
	are equivalent.
	
	By default, the model is estimated by maximum likelihood,
	and the covariance matrix is computed using the negative
	inverse of the Hessian. The \pythcode{opg} option selects
	the Outer Product of the Gradient covariance matrix,
	while \pythcode{robust} requests QML standard errors.
	
	With \pythcode{two\_step=True}, Heckman's two-step estimator
	is used instead. In this case, the covariance matrix for
	the outcome equation is adjusted to account for the
	two-step estimation procedure.}

\gretlfpycmd{hsk()}
\estdesc{Specify a heteroskedasticity-corrected linear regression
	model, estimated by feasible weighted least squares.
	The error variance is modeled as a function of the regressors.}
\begin{lstlisting}[style=pythonstylefoot]
gretl.hsk(formula: str | list[, OPTS]) -> gretl.Model_hsk
\end{lstlisting}
\mapto{hsk depvar indepvars}
\begin{itemize}
	\item[]\maptoopt{no\_squares=True}{--no-squares}{exclude squared regressors from the variance equation}
	\item[]\maptoopt{vcv=True}{--vcv}{print covariance matrix}
	\item[]\maptoopt{quiet=True}{--quiet}{suppress printing of results}
\end{itemize}
\generaldesc{The estimator is designed for models in which
	heteroskedasticity can be approximated by a quadratic
	function of the regressors. It may provide more efficient
	coefficient estimates and consistent standard errors
	compared with OLS.
	
	Estimation proceeds in three stages. First, the model
	is estimated by OLS. Second, an auxiliary regression
	of the logarithm of squared OLS residuals is performed
	on the original regressors and, by default, their squares.
	With \pythcode{no\_squares=True}, the squared regressors
	are omitted.
	
	Finally, weighted least squares is performed using
	the reciprocal of the estimated conditional error variance.
	Specifically, the weights are calculated as
	$w_i = 1/\exp(\hat{u}_i^*)$, where $\hat{u}_i^*$
	denotes the fitted value from the auxiliary regression.
	The exponential transformation ensures that the estimated
	variances are positive.}

\gretlfpycmd{intreg()}
\estdesc{Specify an interval regression model for estimation
	by maximum likelihood, where the dependent variable is
	observed within an interval rather than necessarily
	as an exact value.}
\begin{lstlisting}[style=pythonstylefoot]
gretl.intreg(formula: str | list[, OPTS]) -> gretl.Model_intreg
\end{lstlisting}
\mapto{intreg minvar maxvar indepvars}
\begin{itemize}
	\item[]\maptoopt{quiet=True}{--quiet}{suppress printing of results}
	\item[]\maptoopt{verbose=True}{--verbose}{print details of numerical iterations}
	\item[]\maptoopt{robust=True}{--robust}{use QML robust standard errors}
	\item[]\maptoopt{opg=True}{--opg}{use the outer product of the gradient}
	\item[]\maptoopt{cluster="clustvar"}{--cluster=clustvar}{use clustered standard errors}
\end{itemize}
\noRformula\\\\
\generaldesc{The model assumes a latent dependent variable
	$y_i^* = x_i'\beta + u_i$, where the disturbance term
	is normally distributed. Instead of observing $y_i^*$
	directly, the data provide lower and upper bounds,
	$m_i \leq y_i^* \leq M_i$.
	
	The first two variables in the model specification
	define these bounds. Equal bounds represent exact
	observations. Missing values in \pythcode{minvar}
	indicate left-unbounded intervals, while missing
	values in \pythcode{maxvar} indicate right-unbounded
	intervals.
	
	Estimation is performed by maximum likelihood.
	By default, the covariance matrix is computed as
	the negative inverse of the Hessian. With
	\pythcode{robust=True}, a QML (Huber--White)
	sandwich covariance matrix is used instead.
	Alternatively, \pythcode{opg=True} requests a
	covariance matrix based solely on the outer product
	of the gradient.}

\gretlfpycmd{lad()}
\estdesc{Calculates a regression that minimizes the sum of
	absolute deviations between observed and fitted values
	of the dependent variable.}
\begin{lstlisting}[style=pythonstylefoot]
gretl.lad(formula: str | list[, OPTS]) -> gretl.Model_lad
\end{lstlisting}
\mapto{lad depvar indepvars}
\begin{itemize}
	\item[]\maptoopt{vcv=True}{--vcv}{print covariance matrix}
	\item[]\maptoopt{no\_vcv=True}{--no-vcv}{skip computation of covariance matrix and standard errors}
	\item[]\maptoopt{quiet=True}{--quiet}{suppress printing of results}
\end{itemize}
\generaldesc{Coefficient estimates are obtained using the
	Barrodale--Roberts simplex algorithm. A warning is issued
	if the solution is not unique.
	
	By default, standard errors are computed using 500 bootstrap
	replications. The covariance matrix, printed when
	\pythcode{vcv=True} is specified, is based on the same
	bootstrap procedure. The \pythcode{no\_vcv=True} option
	skips this computation, in which case standard errors
	are not available.
	
	LAD estimation may be relatively slow for large samples
	or models with many regressors. In such cases, median
	regression using \gretlfpycmd{quantreg()} may be preferable,
	since it uses the faster Frisch--Newton algorithm and
	provides analytical rather than bootstrap standard errors.}

\gretlfpycmd{logistic()}
\estdesc{Specify a linear model using a logistic transformation
	of the dependent variable, estimated by ordinary least squares.
	The dependent variable must be strictly positive.}
\begin{lstlisting}[style=pythonstylefoot]
gretl.logistic(formula: str | list[, OPTS]) -> gretl.Model_logistic
\end{lstlisting}
\mapto{logistic depvar indepvars}
\begin{itemize}
	\item[]\maptoopt{ymax=50}{--ymax=50}{specify the asymptotic maximum of the dependent variable}
	\item[]\maptoopt{robust=True}{--robust}{use robust standard errors}
	\item[]\maptoopt{cluster="clustvar"}{--cluster=clustvar}{compute clustered standard errors}
	\item[]\maptoopt{vcv=True}{--vcv}{print covariance matrix}
	\item[]\maptoopt{fixed\_effects=True}{--fixed-effects}{estimate with group fixed effects}
	\item[]\maptoopt{quiet=True}{--quiet}{suppress printing of results}
\end{itemize}
\generaldesc{The model is estimated by applying ordinary least
	squares to a logistic transformation of the dependent variable.
	The dependent variable must be strictly positive and bounded
	above by an asymptotic maximum, denoted by $y^*$.
	
	If all observed values lie between 0 and 1, the default value
	of $y^*$ is 1; otherwise, if they lie between 0 and 100, the
	default is 100. A different maximum can be specified using
	the \pythcode{ymax} option, provided it exceeds all observed
	values of the dependent variable.
	
	Fitted values and residuals are transformed back to the
	original scale using the inverse logistic transformation.
	Since this transformation is nonlinear, the reported values
	involve an approximation.
	
	For panel data, \pythcode{fixed\_effects=True} enables
	estimation with individual fixed effects by subtracting
	group means from the transformed dependent variable.
	
	For a binary dependent variable, the \gretlfpycmd{logit()}
	estimator should be used instead.}

\gretlfpycmd{logit()}
\estdesc{Specify a binary, ordered, or multinomial logit model
	for estimation by maximum likelihood.}
\begin{lstlisting}[style=pythonstylefoot]
gretl.logit(formula: str | list[, OPTS]) -> gretl.Model_logit
\end{lstlisting}
\mapto{logit depvar indepvars}
\begin{itemize}
	\item[]\maptoopt{robust=True}{--robust}{use QML (Huber--White) robust standard errors}
	\item[]\maptoopt{cluster="clustvar"}{--cluster=clustvar}{compute clustered standard errors}
	\item[]\maptoopt{multinomial=True}{--multinomial}{estimate multinomial logit}
	\item[]\maptoopt{vcv=True}{--vcv}{print covariance matrix}
	\item[]\maptoopt{verbose=True}{--verbose}{print details of numerical iterations}
	\item[]\maptoopt{quiet=True}{--quiet}{suppress printing of results}
	\item[]\maptoopt{p\_values=True}{--p-values}{show p-values instead of slopes}
	\item[]\maptoopt{estrella=True}{--estrella}{use Estrella's pseudo-R-squared}
\end{itemize}
\generaldesc{For a binary dependent variable, the model is estimated
	using binary logit. If the dependent variable is non-binary but
	discrete, ordered logit is estimated by default. Multinomial
	logit can be requested using \pythcode{multinomial=True}.
	
	For binary logit, the estimated coefficients represent effects
	on the log odds. By default, slopes evaluated at the means
	of the regressors are reported instead of p-values.
	This behavior can be disabled using \pythcode{p\_values=True}.
	
	Standard errors are based on the negative inverse of the Hessian
	by default. The \pythcode{robust} option requests QML
	(Huber--White) standard errors, while \pythcode{cluster}
	requests cluster-robust standard errors.
	
	By default, McFadden's pseudo-R-squared is reported.
	For binary logit, Estrella's alternative can be selected
	using \pythcode{estrella=True}.}

\gretlfpycmd{midasreg()}
\estdesc{Specify a Mixed Data Sampling (MIDAS) regression model
	for least-squares estimation, combining regressors observed
	at different frequencies. Depending on the parameterization,
	estimation uses OLS or nonlinear least squares (NLS).}
\begin{lstlisting}[style=pythonstylefoot]
gretl.midasreg(formula: str[, OPTS]) -> gretl.Model_midasreg
\end{lstlisting}
\mapto{midasreg depvar indepvars ; MIDAS-terms}
\begin{itemize}
	\item[]\maptoopt{vcv=True}{--vcv}{print covariance matrix}
	\item[]\maptoopt{robust=True}{--robust}{use robust standard errors}
	\item[]\maptoopt{quiet=True}{--quiet}{suppress printing of results}
	\item[]\maptoopt{levenberg=True}{--levenberg}{use the Levenberg--Marquardt algorithm}
\end{itemize}
\generaldesc{The model specification consists of the dependent
	variable and regressors observed at the same frequency,
	followed by a semicolon and one or more higher-frequency
	MIDAS terms. The latter are specified using either
	\pythcode{mds()} or \pythcode{mdsl()}.
	
	The \pythcode{mds()} specification takes a MIDAS list,
	minimum and maximum lag orders, a parameterization type,
	and an optional initialization argument. The
	\pythcode{mdsl()} specification uses a list containing
	the required lags, followed by the parameterization type
	and an optional initialization argument.
	
	Six parameterizations are supported:
	\pythcode{umidas} (0), unrestricted MIDAS;
	\pythcode{nealmon} (1), normalized exponential Almon;
	\pythcode{beta0} (2), normalized beta with zero last lag;
	\pythcode{betan} (3), normalized beta with non-zero last lag;
	\pythcode{almonp} (4), non-normalized Almon polynomial;
	and \pythcode{beta1} (5), beta with the first parameter
	fixed at one.
	
	U-MIDAS specifications are estimated by OLS, while the
	remaining parameterizations require nonlinear least squares.
	Initial values for nonlinear parameters may be supplied
	using a vector, requested automatically using
	\pythcode{null} where supported, or specified by giving
	the required number of parameters.
	
	For normalized exponential Almon and normalized beta
	parameterizations, the default estimation procedure
	combines constrained BFGS and OLS. The
	\pythcode{levenberg=True} option forces use of the
	Levenberg--Marquardt algorithm.}

\gretlfpycmd{mpols()}
\estdesc{Specify a linear regression model for high-precision
	OLS estimation using multiple-precision floating-point
	arithmetic provided by the GNU Multiple Precision (GMP)
	library.}
\begin{lstlisting}[style=pythonstylefoot]
gretl.mpols(formula: str | list[, OPTS]) -> gretl.Model_mpols
\end{lstlisting}
\mapto{mpols depvar indepvars [; powers]}
\begin{itemize}
	\item[]\maptoopt{vcv=True}{--vcv}{print covariance matrix}
	\item[]\maptoopt{simple\_print=True}{--simple-print}{suppress auxiliary model statistics}
	\item[]\maptoopt{quiet=True}{--quiet}{suppress printing of results}
\end{itemize}
\generaldesc{Estimation uses multiple-precision floating-point
	arithmetic, with 256 bits of precision by default.
	The precision can be increased by setting the
	\pythcode{GRETL\_MP\_BITS} environment variable before
	starting Python.
	
	Optionally, a list of powers may be specified after a
	semicolon. These powers are applied to the last regressor
	in the model specification, and the resulting terms are
	computed and stored using multiple-precision arithmetic.
	
	When using Python lists, the semicolon-separated components
	are represented by nested lists. For example, the specifications
	\pythcode{"y 0 x ; 2 3 4"} and
	\pythcode{[["y", 0, "x"], [2, 3, 4]]}
	are equivalent. Both specify a regression of $y$ on
	a constant, $x$, $x^2$, $x^3$ and $x^4$.}

\gretlfpycmd{negbin()}
\estdesc{Specify a Negative Binomial regression model for estimation
	by maximum likelihood. The dependent variable represents a count
	and must contain only non-negative integer values.}
\begin{lstlisting}[style=pythonstylefoot]
gretl.negbin(formula: str | list[, OPTS]) -> gretl.Model_negbin
\end{lstlisting}
\mapto{negbin depvar indepvars [; offset]}
\begin{itemize}
	\item[]\maptoopt{model1=True}{--model1}{use the NegBin 1 specification}
	\item[]\maptoopt{robust=True}{--robust}{use QML robust standard errors}
	\item[]\maptoopt{cluster="clustvar"}{--cluster=clustvar}{compute clustered standard errors}
	\item[]\maptoopt{opg=True}{--opg}{use the Outer Product of the Gradient covariance matrix}
	\item[]\maptoopt{vcv=True}{--vcv}{print covariance matrix}
	\item[]\maptoopt{verbose=True}{--verbose}{print details of numerical iterations}
	\item[]\maptoopt{quiet=True}{--quiet}{suppress printing of results}
\end{itemize}
\generaldesc{By default, the NegBin~2 specification is estimated,
	with conditional variance $\mu(1+\alpha\mu)$, where $\mu$
	denotes the conditional mean. The \pythcode{model1=True}
	option selects the NegBin 1 specification, with conditional
	variance $\mu(1+\alpha)$. The Poisson model corresponds
	to the restricted case $\alpha=0$.
	
	An optional offset variable may be specified after a semicolon
	in the model specification. Its logarithm enters the linear
	predictor with a coefficient fixed at 1. The offset variable
	must be strictly positive.
	
	When using Python lists, semicolon-separated variable lists
	are represented by nested lists. For example, the specifications
	\pythcode{"y 0 x1 x2 ; exposure"} and
	\pythcode{[["y", "const", "x1", "x2"], ["exposure"]]}
	are equivalent.
	
	By default, standard errors are based on a numerical
	approximation to the Hessian. The \pythcode{opg} option
	selects the Outer Product of the Gradient covariance matrix,
	while \pythcode{robust} requests QML sandwich standard errors.
	The \pythcode{cluster} option computes cluster-robust
	standard errors.}

\gretlfpycmd{nls()}
\estdesc{Performs Nonlinear Least Squares (NLS) estimation using
	a modified version of the Levenberg-Marquardt algorithm -- \experimental.}
\begin{lstlisting}[style=pythonstylefoot]
gretl.nls(formula: str[, data: gretl.Dataset, OPTS]) -> gretl.Model_nls
\end{lstlisting}
\mapto[none]{nls function [ derivatives ]}
\generaldesc{\tbw}

\gretlfpycmd{ols()}
\estdesc{Specify a linear model for estimation by ordinary least squares (OLS).}
\begin{lstlisting}[style=pythonstylefoot]
gretl.ols(formula: str | list[, OPTS]) -> gretl.Model_ols
\end{lstlisting}
\mapto{ols depvar indepvars}
\begin{itemize}
	\item[]\maptoopt{vcv=True}{--vcv}{print covariance matrix}
	\item[]\maptoopt{robust=True}{--robust}{use heteroskedasticity-robust standard errors}
	\item[]\maptoopt{cluster="clustvar"}{--cluster=clustvar}{compute clustered standard errors}
	\item[]\maptoopt{jackknife=True}{--jackknife}{use the HC3a variant of robust standard errors}
	\item[]\maptoopt{simple\_print=True}{--simple-print}{suppress auxiliary model statistics}
	\item[]\maptoopt{quiet=True}{--quiet}{suppress printing of results}
	\item[]\maptoopt{anova=True}{--anova}{print an ANOVA table}
	\item[]\maptoopt{no\_df\_corr=True}{--no-df-corr}{suppress degrees-of-freedom correction}
	\item[]\maptoopt{print\_final=True}{--print-final}{print results only on the final loop iteration}
\end{itemize}
\generaldesc{The model may be specified using a \lgretl-style
	variable list, a formula, or a Python list of variable names
	or IDs. The dependent variable comes first, and the value
	0 denotes the constant term.
	
	The \pythcode{robust} option computes heteroskedasticity-robust
	standard errors using the currently selected HC estimator.
	The \pythcode{jackknife} option selects HC3a, while
	\pythcode{cluster} produces standard errors clustered by
	the distinct values of the specified variable.
	
	By default, the estimated error variance uses a
	degrees-of-freedom correction, which can be disabled using
	\pythcode{no\_df\_corr=True}.
	
	The \pythcode{print\_final} option is applicable only within
	loops and suppresses estimation output until the final iteration.}

\gretlfpycmd{panel()}
\estdesc{Specify a linear panel-data model for estimation.
	By default, the fixed-effects estimator is used, but
	random-effects, pooled OLS, between-groups and
	unit-weighted estimators are also available.}
\begin{lstlisting}[style=pythonstylefoot]
gretl.panel(formula: str | list[, OPTS]) -> gretl.Model_panel
\end{lstlisting}
\mapto{panel depvar indepvars}
\begin{itemize}
	\item[]\maptoopt{vcv=True}{--vcv}{print covariance matrix}
	\item[]\maptoopt{fixed\_effects=True}{--fixed-effects}{estimate with group fixed effects}
	\item[]\maptoopt{random\_effects=True}{--random-effects}{estimate a random-effects model}
	\item[]\maptoopt{nerlove=True}{--nerlove}{use the Nerlove method}
	\item[]\maptoopt{pooled=True}{--pooled}{estimate via pooled OLS}
	\item[]\maptoopt{between=True}{--between}{estimate the between-groups model}
	\item[]\maptoopt{robust=True}{--robust}{use robust standard errors \cite[see][]{gretl_command}}
	\item[]\maptoopt{cluster="cvar"}{--cluster=cvar}{compute clustered standard errors}
	\item[]\maptoopt{time\_dummies=True}{--time-dummies}{include time dummy variables}
	\item[]\maptoopt{unit\_weights=True}{--unit-weights}{estimate using unit-weighted least squares}
	\item[]\maptoopt{iterate=True}{--iterate}{use iterative estimation}
	\item[]\maptoopt{matrix\_diff=True}{--matrix-diff}{compute Hausman test via matrix difference}
	\item[]\maptoopt{unbalanced="bc"}{--unbalanced=bc}{use the Baltagi--Chang method for unbalanced panels}
	\item[]\maptoopt{unbalanced="stata"}{--unbalanced=stata}{use the Stata-compatible method for unbalanced panels}
	\item[]\maptoopt{quiet=True}{--quiet}{reduce printed output}
	\item[]\maptoopt{verbose=True}{--verbose}{print additional estimation details}
\end{itemize}
\generaldesc{By default, the model is estimated using group
	fixed effects, implemented by subtracting unit means
	from the original data.
	
	With \pythcode{random\_effects=True}, the Swamy--Arora
	estimator is used by default. The \pythcode{nerlove=True}
	option selects the Nerlove method instead. For random-effects
	models, \pythcode{matrix\_diff=True} requests the
	matrix-difference form of the Hausman test.
	
	Alternatively, \pythcode{pooled=True} estimates pooled OLS,
	while \pythcode{between=True} estimates a regression using
	group means. With \pythcode{unit\_weights=True}, weighted
	least squares is used, with weights based on unit-specific
	residual variances. The \pythcode{iterate=True} option
	enables iterative estimation in this case; convergence
	yields maximum likelihood estimates.
	
	With \pythcode{robust=True}, the default covariance
	estimator is the Arellano HAC estimator, clustered by
	panel unit. Alternative estimators include Panel Corrected
	Standard Errors (PCSE) and Driscoll--Kraay standard errors,
	selectable through gretl's \gretlcmd{set panel_robust}
	command. The \pythcode{cluster} option provides additional
	clustering possibilities.
	
	For unbalanced random-effects models, the
	\pythcode{unbalanced} option selects the ANOVA method.
	Supported values are \pythcode{"bc"} for the
	Baltagi--Chang method and \pythcode{"stata"} for
	compatibility with Stata's \pythcode{xtreg, sa} estimator.}

\gretlfpycmd{poisson()}
\estdesc{Specify a Poisson regression model for estimation
	by maximum likelihood. The dependent variable represents
	a count and must contain only non-negative integer values.}
\begin{lstlisting}[style=pythonstylefoot]
gretl.poisson(formula: str | list[, OPTS]) -> gretl.Model_poisson
\end{lstlisting}
\mapto{poisson depvar indepvars [; offset]}
\begin{itemize}
	\item[]\maptoopt{robust=True}{--robust}{use QML (Huber--White) robust standard errors}
	\item[]\maptoopt{cluster="clustvar"}{--cluster=clustvar}{compute clustered standard errors}
	\item[]\maptoopt{vcv=True}{--vcv}{print covariance matrix}
	\item[]\maptoopt{verbose=True}{--verbose}{print details of numerical iterations}
	\item[]\maptoopt{quiet=True}{--quiet}{suppress printing of results}
\end{itemize}
\noRformula\\
\generaldesc{The conditional mean of the dependent variable is modeled
	as an exponential function of the regressors, corresponding to
	a logarithmic link function.
	
	An optional offset variable may be specified after a semicolon
	in the model specification. Its logarithm enters the linear
	predictor with a coefficient fixed at 1. The offset variable
	must be strictly positive.
	
	When using Python lists, semicolon-separated variable lists
	are represented by nested lists. For example, the specifications
	\pythcode{"y 0 x1 x2 ; exposure"} and
	\pythcode{[["y", "const", "x1", "x2"], ["exposure"]]}
	are equivalent.
	
	By default, standard errors are computed using the negative
	inverse of the Hessian. The \pythcode{robust} option requests
	QML (Huber--White) standard errors, while \pythcode{cluster}
	requests cluster-robust standard errors.}

\gretlfpycmd{probit()}
\estdesc{Specify a binary, ordered, or random-effects panel probit
	model for estimation by maximum likelihood.}
\begin{lstlisting}[style=pythonstylefoot]
gretl.probit(formula: str | list[, OPTS]) -> gretl.Model_probit
\end{lstlisting}
\mapto{probit depvar indepvars}
\begin{itemize}
	\item[]\maptoopt{robust=True}{--robust}{use QML (Huber--White) robust standard errors}
	\item[]\maptoopt{cluster="clustvar"}{--cluster=clustvar}{compute clustered standard errors}
	\item[]\maptoopt{vcv=True}{--vcv}{print covariance matrix}
	\item[]\maptoopt{verbose=True}{--verbose}{print details of numerical iterations}
	\item[]\maptoopt{quiet=True}{--quiet}{suppress printing of results}
	\item[]\maptoopt{p\_values=True}{--p-values}{show p-values instead of slopes}
	\item[]\maptoopt{estrella=True}{--estrella}{use Estrella's pseudo-R-squared}
	\item[]\maptoopt{random\_effects=True}{--random-effects}{estimate a random-effects panel probit model}
	\item[]\maptoopt{quadpoints=k}{--quadpoints=k}{set the number of Gauss--Hermite quadrature points}
\end{itemize}
\generaldesc{For a binary dependent variable, binary probit is estimated.
	If the dependent variable is non-binary but discrete, ordered probit
	is estimated instead.
	
	For binary probit, slopes evaluated at the means of the regressors
	are reported instead of p-values by default. This behavior can be
	disabled using \pythcode{p\_values=True}.
	
	Standard errors are based on the negative inverse of the Hessian
	by default. The \pythcode{robust} option requests QML
	(Huber--White) standard errors, while \pythcode{cluster}
	requests cluster-robust standard errors.
	
	The \pythcode{random\_effects=True} option enables random-effects
	probit estimation for panel data. The likelihood is evaluated using
	Gauss--Hermite quadrature, with 32 quadrature points by default.
	The number of points can be adjusted using \pythcode{quadpoints}.
	
	By default, McFadden's pseudo-R-squared is reported.
	For binary probit, Estrella's alternative can be selected
	using \pythcode{estrella=True}.}

\gretlfpycmd{system()}
\estdesc{System of equations -- \experimental.}
\begin{lstlisting}[style=pythonstylefoot]
gretl.system(method: str[,data: gretl.Dataset, OPTS]) -> gretl.Model_NME
\end{lstlisting}
\mapto[none]{system method=estimator}
\generaldesc{\tbw}

\gretlfpycmd{quantreg()}
\estdesc{Quantile regression}
\begin{lstlisting}[style=pythonstylefoot]
gretl.quantreg(tau: float | numpy.ndarray, formula: str | list[, OPTS]) -> gretl.Model_quantreg
\end{lstlisting}
\mapto{quantreg tau depvar indepvars}
\begin{itemize}
	\item[]\maptoopt{robust=True}{--robust}{heteroskedasticity-robust standard errors or confidence intervals}
	\item[]\maptoopt{intervals=True}{--intervals}{compute 90\% confidence intervals}
	\item[]\maptoopt{intervals=0.95}{--intervals=0.95}{specify the confidence level}
	\item[]\maptoopt{vcv=True}{--vcv}{print covariance matrix}
	\item[]\maptoopt{quiet=True}{--quiet}{suppress printing of results}
\end{itemize}
\generaldesc{The argument \textit{tau} specifies a conditional
	quantile or a sequence of quantiles, with values between
	0.01 and 0.99. The regression specification follows the
	same conventions as \gretlfpycmd{ols()}.
	
	By default, standard errors are computed using the asymptotic
	method of Koenker and Bassett (1978). The \pythcode{robust}
	option selects the heteroskedasticity-robust method of
	Koenker and Zhao (1994).
	
	The \textit{intervals} option produces confidence intervals
	instead of standard errors, using rank inversion. The default
	confidence level is 90\%. With \pythcode{robust=True}, the
	intervals follow Koenker and Machado (1999).
	
	When multiple quantiles are specified, the coefficient and
	standard error results are matrices with one column per
	quantile. The combined results are also available under
	\pythcode{rq_sequence} in the model results.}

\gretlfpycmd{tobit()}
\estdesc{Specify a Tobit model for maximum-likelihood estimation.}
\begin{lstlisting}[style=pythonstylefoot]
gretl.tobit(formula: str | list[, OPTS]) -> gretl.Model_tobit
\end{lstlisting}
\mapto{tobit depvar indepvars}
\begin{itemize}
	\item[]\maptoopt{llimit=lval}{--llimit=lval}{specify the left censoring bound}
	\item[]\maptoopt{rlimit=rval}{--rlimit=rval}{specify the right censoring bound}
	\item[]\maptoopt{vcv=True}{--vcv}{print covariance matrix}
	\item[]\maptoopt{robust=True}{--robust}{use QML robust standard errors}
	\item[]\maptoopt{opg=True}{--opg}{use the outer product of the gradient for the covariance matrix}
	\item[]\maptoopt{cluster="clustvar"}{--cluster=clustvar}{compute cluster-robust standard errors}
	\item[]\maptoopt{verbose=True}{--verbose}{print details of numerical iterations}
	\item[]\maptoopt{quiet=True}{--quiet}{suppress printing of results}
\end{itemize}
\generaldesc{The Tobit model is appropriate for a censored dependent
	variable. By default, the dependent variable is assumed to be
	left-censored at zero and uncensored on the right.
	
	The \textit{llimit} and \textit{rlimit} options specify
	alternative censoring bounds. If only \pythcode{rlimit} is
	specified, the dependent variable is assumed to be uncensored
	on the left.
	
	By default, the covariance matrix is based on the negative
	inverse Hessian. The \pythcode{robust} option selects QML
	(Huber--White) robust standard errors, while \pythcode{opg}
	uses the outer product of the gradient. The \pythcode{cluster}
	option specifies a variable identifying clusters for
	cluster-robust inference.
	
	The Tobit model is a special case of interval regression.
	See also \gretlfpycmd{intreg()}.}

\gretlfpycmd{tsls()}
\estdesc{Specify an instrumental-variables model for estimation,
	by default using two-stage least squares (TSLS).
	Limited Information Maximum Likelihood (LIML) and
	Generalized Method of Moments (GMM) estimation are
	also supported.}
\begin{lstlisting}[style=pythonstylefoot]
gretl.tsls(formula: str | list[, OPTS]) -> gretl.Model_tsls
\end{lstlisting}
\mapto{tsls depvar indepvars ; instruments}
\begin{itemize}
	\item[]\maptoopt{no\_tests=True}{--no-tests}{suppress diagnostic tests}
	\item[]\maptoopt{vcv=True}{--vcv}{print covariance matrix}
	\item[]\maptoopt{quiet=True}{--quiet}{suppress printing of results}
	\item[]\maptoopt{no\_df\_corr=True}{--no-df-corr}{suppress degrees-of-freedom correction}
	\item[]\maptoopt{robust=True}{--robust}{use robust standard errors}
	\item[]\maptoopt{cluster="clustvar"}{--cluster=clustvar}{use clustered standard errors}
	\item[]\maptoopt{matrix\_diff=True}{--matrix-diff}{use matrix-difference Hausman test}
	\item[]\maptoopt{liml=True}{--liml}{use Limited Information Maximum Likelihood}
	\item[]\maptoopt{gmm=True}{--gmm}{use the Generalized Method of Moments}
	\item[]\maptoopt{two\_step=True}{--two-step}{use two-step GMM}
	\item[]\maptoopt{iterate=True}{--iterate}{iterate GMM to convergence}
	\item[]\maptoopt{weights="Wmat"}{--weights=Wmat}{specify GMM weighting matrix}
\end{itemize}
\generaldesc{The model specification consists of two variable
	lists separated by a semicolon. The first specifies the
	dependent variable and regressors, while the second
	specifies the instruments. The instrument list must
	contain at least as many variables as the regressor
	list. Exogenous regressors should appear in both lists.
	
	When using Python lists, the two components are represented
	by nested lists. For example, the specifications
	\pythcode{"y 0 x1 x2 ; 0 x1 z1 z2"} and
	\pythcode{[["y", 0, "x1", "x2"], [0, "x1", "z1", "z2"]]}
	are equivalent.
	
	By default, TSLS estimation includes the Hausman test
	of OLS consistency and, for overidentified models,
	the Sargan test of overidentifying restrictions.
	The corresponding statistics are available through
	the \gretlcmd{\$hausman} and \gretlcmd{\$sargan}
	accessors, respectively.
	
	For TSLS and LIML estimation under i.i.d.\ errors,
	a weak-instruments test is also reported. With a single
	endogenous regressor, the test is based on the first-stage
	F statistic; otherwise, it uses the smallest eigenvalue
	of the corresponding matrix statistic.
	
	The \pythcode{matrix\_diff=True} option selects the
	matrix-difference method for the Hausman test instead
	of the default regression-based method.
	
	Alternatively, \pythcode{liml=True} selects Limited
	Information Maximum Likelihood estimation, while
	\pythcode{gmm=True} selects the Generalized Method
	of Moments. For just-identified models, these estimators
	yield equivalent results, but they generally differ
	for overidentified models.
	
	For GMM estimation, \pythcode{two\_step=True} enables
	two-step estimation, \pythcode{iterate=True} iterates
	the estimator to convergence, and \pythcode{weights}
	specifies a weighting matrix. By default, the GMM
	weighting matrix is the identity matrix.}

\gretlfpycmd{wls()}
\estdesc{Specify a linear model for estimation by weighted least squares (WLS).}
\begin{lstlisting}[style=pythonstylefoot]
	gretl.wls(formula: str | list[, OPTS]) -> gretl.Model_wls
\end{lstlisting}
\mapto{wls wtvar depvar indepvars}
\begin{itemize}
	\item[]\maptoopt{vcv=True}{--vcv}{print covariance matrix}
	\item[]\maptoopt{robust=True}{--robust}{use heteroskedasticity-robust standard errors}
	\item[]\maptoopt{quiet=True}{--quiet}{suppress printing of results}
	\item[]\maptoopt{allow\_zeros=True}{--allow-zeros}{allow zero weights together with positive weights}
\end{itemize}
\noRformula\\
\generaldesc{The first variable specifies the weights, followed by
	the dependent variable and regressors. The model can be specified
	using a \lgretl-style variable list or a Python list of variable
	names or IDs.
	
	WLS estimation is equivalent to OLS applied to the variables
	multiplied by the square root of the weights, although the
	coefficient of determination is computed differently.
	
	The weighting variable should normally contain strictly positive
	values. If it is a binary variable, WLS is equivalent to OLS
	on the subsample with nonzero weights. Other combinations of
	zero and positive weights require \pythcode{allow\_zeros=True}.}

\gretlfpycmd{var()}
\estdesc{Specify a vector autoregression (VAR) model for estimation
	by ordinary least squares.}
\begin{lstlisting}[style=pythonstylefoot]
gretl.var(order: int, ylist: str | list[, OPTS]) -> gretl.Model_var
\end{lstlisting}
\mapto{var order ylist [; xlist]}
\begin{itemize}
	\item[]\maptoopt{xlist="x1 x2"}{; xlist}{include exogenous regressors}
	\item[]\maptoopt{nc=True}{--nc}{do not include a constant}
	\item[]\maptoopt{trend=True}{--trend}{include a linear trend}
	\item[]\maptoopt{seasonals=True}{--seasonals}{include seasonal dummy variables}
	\item[]\maptoopt{robust=True}{--robust}{heteroskedasticity-robust standard errors}
	\item[]\maptoopt{robust\_hac=True}{--robust-hac}{HAC standard errors}
	\item[]\maptoopt{quiet=True}{--quiet}{suppress output of individual equations}
	\item[]\maptoopt{silent=True}{--silent}{suppress all printed output}
	\item[]\maptoopt{impulse\_responses=True}{--impulse-responses}{print impulse responses}
	\item[]\maptoopt{variance\_decomp=True}{--variance-decomp}{print forecast-error variance decompositions}
	\item[]\maptoopt{lagselect=True}{--lagselect}{perform lag-order selection}
	\item[]\maptoopt{minlag=0}{--minlag=0}{specify minimum lag order for lag selection}
\end{itemize}
\noRformula\\
\generaldesc{The argument \pythcode{order} specifies the VAR lag
	order, while \pythcode{ylist} identifies the endogenous variables.
	Their lags are generated automatically. Optional exogenous
	regressors can be specified via \pythcode{xlist}.
	
	A constant is included by default, unless \pythcode{nc=True}
	is specified. A linear trend and seasonal dummy variables may
	be added using the corresponding options.
	
	Impulse responses and forecast-error variance decompositions
	are based on the Cholesky decomposition of the contemporaneous
	covariance matrix. Their results therefore depend on the
	ordering of endogenous variables in \pythcode{ylist}.
	
	With \pythcode{lagselect=True}, \pythcode{order} is interpreted
	as the maximum lag order. The resulting table reports
	likelihood-ratio tests and the AIC, BIC and HQC information
	criteria. The \pythcode{minlag} option controls the minimum
	lag order considered.}

\gretlfpycmd{vecm()}
\estdesc{Specify a vector error-correction model (VECM) for
	cointegrated time series.}
\begin{lstlisting}[style=pythonstylefoot]
gretl.vecm(order: int, rank: int, ylist: str | list[, OPTS]) -> gretl.Model_vecm
\end{lstlisting}
\mapto{vecm order rank ylist [; xlist] [; rxlist]}
\begin{itemize}
	\item[]\maptoopt{nc=True}{--nc}{include no constant}
	\item[]\maptoopt{rc=True}{--rc}{include a restricted constant}
	\item[]\maptoopt{uc=True}{--uc}{include an unrestricted constant}
	\item[]\maptoopt{crt=True}{--crt}{include a constant and restricted trend}
	\item[]\maptoopt{ct=True}{--ct}{include a constant and unrestricted trend}
	\item[]\maptoopt{seasonals=True}{--seasonals}{include centered seasonal dummy variables}
	\item[]\maptoopt{quiet=True}{--quiet}{suppress output of individual equations}
	\item[]\maptoopt{silent=True}{--silent}{suppress all printed output}
	\item[]\maptoopt{impulse\_responses=True}{--impulse-responses}{print impulse responses}
	\item[]\maptoopt{variance\_decomp=True}{--variance-decomp}{print forecast-error variance decompositions}
\end{itemize}
\noRformula\\
\generaldesc{The argument \pythcode{order} specifies the lag order
	of the underlying VAR system. The corresponding VECM contains
	\pythcode{order - 1} lags of the differenced endogenous variables.
	The argument \pythcode{rank} specifies the cointegration rank,
	which must be positive and cannot exceed the number of
	endogenous variables.
	
	The variables in \pythcode{ylist} should be integrated of
	order one and are specified in levels. Optional exogenous
	variables may be supplied through \pythcode{xlist} and
	\pythcode{rxlist}, separated by semicolons. Variables in
	\pythcode{xlist} enter the model unrestrictedly, while those
	in \pythcode{rxlist} are restricted to the cointegration space.
	
	By default, an unrestricted constant is included (Johansen's
	case 3). The mutually exclusive options \pythcode{nc},
	\pythcode{rc}, \pythcode{crt} and \pythcode{ct} select
	alternative deterministic specifications (cases 1, 2, 4
	and 5, respectively). The \pythcode{uc} option explicitly
	selects the default specification.
	
	The \pythcode{seasonals} option adds centered seasonal dummy
	variables and is available for quarterly and monthly data.
	
	Impulse responses and forecast-error variance decompositions
	are based on Cholesky identification and therefore depend
	on the ordering of endogenous variables.
	
	Following estimation, the Johansen adjustment and
	cointegration matrices are available through the
	\pythcode{jalpha} and \pythcode{jbeta} model results,
	while \pythcode{jvbeta} provides the estimated variance
	matrix of the cointegrating coefficients.}
